\section{Conclusions}

Functional intent in T3SD is classically discharged by exhibiting a homomorphism from a build onto a reference system. We identify a temporal fragment---readout and next-state clauses compiled from the reference table, with no eventually and no until---for which satisfaction is equivalent to existence of such a homomorphism, and we argue that engineers should treat the fragment as the artifact they maintain. The reason is asymmetry rather than novelty: the fragment depends on the reference alone, so it is invariant under state re-encoding, port renaming and permutation, architectural regrouping, and zone-by-zone rebuilding, and it pins the reference down to isomorphism among finite systems. A homomorphism is data about a particular pair of encodings and must be rebuilt whenever either changes.

The case study is a Turing machine presented as a Wymore coupling: a tape zone and a two-phase control zone with declared ports, a two-wire connectivity, and a resultant shown to realize a monolithic reference. Universality is definitional for that construction, so the demonstration covers arbitrary computable functional intent rather than a convenient example. Its zones are then rebuilt in several ways---accepted variants and defective ones---and conformance of the rebuilt machine follows from conformance of its zones without the machine being re-examined. All of these results, including the rejections, are machine-checked.

Because satisfaction of the fragment is a per-state condition, the conformance question is mechanically decidable over finite systems, and we report a checker that decides it: on a finite instance suite it returns a conformance map where one exists and a proof of non-existence where none does, every verdict agrees with its machine-checked counterpart, and one solver-found map is re-verified by the proof assistant. Two honest limits bound that: existence of a surjective conformance map is NP-hard, so the timings are feasibility evidence rather than a scalability claim, and what makes larger systems tractable in this setting is the composition result rather than solver capacity.

Two limitations point at the work we think is most worth doing next. Conformance is lock-step, so a build that takes two ticks to accomplish what intent states in one is rejected---we prove this---and intent must be restated at the build's granularity; a rate-tolerant conformance relation would be a weaker relation than Wymore's and needs its own theory. And the fragment cannot express liveness, so termination and response intent must be argued in a richer dialect, outside the bridge. Global equivalence for properties beyond the fragment remains open, as do toolchain integration and measurement of engineering effort.
